% =====================================================================
% Capítulo 3. Marco teórico y estado del arte
% Estado: 3.1, 3.2 y 3.3 redactadas (06/10/2026); 3.4 a 3.8 pendientes.
% =====================================================================

\chapter{Marco teórico y estado del arte}\label{cap:marco-teorico}

Este capítulo expone los fundamentos sobre los que se apoya el trabajo. Se
organiza en ocho apartados. Los tres primeros tratan la demanda eléctrica
como serie temporal, los fundamentos de los modelos clásicos de series
temporales y su extensión a series con varias estacionalidades. Los
siguientes presentan los métodos de aprendizaje automático y profundo
aplicados a la previsión, los criterios de evaluación de previsiones
puntuales, la previsión probabilística y su calibración, y, por último, el
estado del arte en previsión de la demanda eléctrica, del que surgen las
carencias que motivan el trabajo. El capítulo explica en qué consiste cada
método; las decisiones concretas sobre cómo se aplican se recogen en el
marco metodológico.

% =====================================================================
\section{La demanda eléctrica como serie temporal}\label{sec:demanda-serie}

La sección~\ref{sec:contexto} presentó el papel de la previsión de la demanda
en la operación del sistema y en el mercado. Esta sección se ocupa de la
demanda como objeto de modelización: qué se mide, de qué depende y qué tipos
de previsión pueden construirse sobre ella.

\subsection{Qué se mide}\label{sec:que-se-mide}

La demanda eléctrica es la potencia que los consumidores toman de la red en
cada instante, y se expresa en megavatios (MW). Su integral a lo largo de un
intervalo es la energía consumida, expresada en megavatios hora (MWh). Las
series de demanda se publican como valores muestreados cada pocos minutos o
como valores medios por intervalo; la media de la potencia a lo largo de una
hora coincide numéricamente con la energía consumida en esa hora, de modo
que en una serie horaria ambas lecturas son equivalentes.

La demanda que observa el operador del sistema es la que debe cubrirse con
la generación conectada a la red. Es el agregado de millones de consumos
individuales, domésticos, comerciales e industriales, cada uno de ellos muy
irregular. La agregación compensa buena parte de esa irregularidad, y el
resultado es una serie suave y con regularidades muy marcadas, lo que
explica que su previsión a corto plazo pueda alcanzar errores relativos de
pocas unidades porcentuales. Esa misma agregación implica que la serie no
distingue entre tipos de consumidores, y que los cambios en la composición
del consumo, como la expansión del autoconsumo, aparecen en ella como
cambios de nivel o de forma sin una variable que los explique.

\subsection{Determinantes de la demanda}\label{sec:determinantes}

Los factores que influyen en la demanda a corto plazo pueden agruparse en
cuatro tipos.

El primero es el ritmo de la actividad humana, que genera los ciclos
regulares de la serie. El ciclo diario refleja la sucesión de un valle
nocturno, una subida por la mañana con el inicio de la actividad y los
máximos de mediodía y de última hora de la tarde. El ciclo semanal distingue
los días laborables, de perfil similar entre sí, del sábado y, sobre todo,
del domingo, en los que la actividad industrial y comercial se reduce. Como
la actividad sigue el reloj oficial, estos ciclos se organizan según la hora
local y no según el tiempo universal.

El segundo es el calendario. Los días festivos se comportan de forma
parecida a un domingo, y los puentes, días laborables situados entre un
festivo y un fin de semana, presentan una actividad intermedia. En un país
con festivos regionales, como España, un festivo autonómico reduce la
demanda solo en la parte del territorio afectada. Hay además periodos sin
carácter festivo pero con actividad reducida, como el periodo navideño. Los
efectos de calendario son difíciles de modelizar porque cada tipo de día
especial aparece pocas veces al año y su efecto puede variar según el día de
la semana en que caiga.

El tercero son las condiciones meteorológicas, y en particular la
temperatura. La demanda aumenta con el frío, por el uso de la calefacción
eléctrica, y con el calor, por el de la refrigeración, de modo que la
relación con la temperatura es no lineal, con una zona intermedia de baja
sensibilidad. Esta forma ha sido documentada para la demanda española
\parencite{valor2001daily,moralcarcedo2005modelling}. Una forma habitual de
representarla son los grados-día. Dado un umbral de temperatura $T^{*}$, los
grados-día de calefacción de un día con temperatura media $T$ son
$\max(T^{*} - T, 0)$ y los de refrigeración, $\max(T - T^{*}, 0)$, de modo
que cada variable solo es positiva en su rama y crece linealmente al
alejarse del umbral. Se pueden emplear dos umbrales distintos para delimitar
una zona neutra en la que ninguna de las dos se activa. La elección de los
umbrales depende del clima y de los equipamientos de cada sistema, por lo que
conviene estimarlos con los datos. Otras variables meteorológicas, como la
humedad, la nubosidad o la radiación solar, también influyen, la última de
forma creciente por su efecto sobre el autoconsumo fotovoltaico, aunque su
peso en la demanda agregada suele ser menor que el de la temperatura.

El cuarto son los factores estructurales: la actividad económica, la
eficiencia de los equipos, la electrificación de nuevos usos, la difusión
del autoconsumo o los precios. Evolucionan despacio y se manifiestan en la
serie como una tendencia o como cambios de nivel entre años. Algunos
episodios excepcionales, como las restricciones de movilidad de la pandemia
de 2020, producen además roturas bruscas que interrumpen temporalmente la
relación entre la demanda y sus determinantes habituales.

\subsection{Horizontes y tipos de previsión}\label{sec:tipos-prevision}

La previsión de la demanda se realiza a horizontes muy distintos, desde
minutos hasta años, y cada horizonte da protagonismo a determinantes
diferentes. A pocas horas vista, la demanda reciente es el mejor indicador
de la próxima. A uno o varios días, dominan el calendario y la temperatura
prevista, que determinan la forma del perfil diario. A meses o años, pierden
peso los factores meteorológicos concretos y lo ganan los estructurales y la
climatología media. Este trabajo se sitúa en el corto plazo, con un horizonte
de entre 12 y 36 horas (sección~\ref{sec:problema}), en el que la
temperatura y el calendario son los determinantes principales.

Una previsión emitida en un momento dado para varias horas futuras es una
previsión multihorizonte: el modelo debe producir de una vez los 24 valores
del día siguiente, y cada uno de ellos tiene una anticipación distinta. Las
estrategias para hacerlo, recursiva o directa, se tratan en la
sección dedicada al aprendizaje profundo.

Por último, una previsión puede ser puntual o probabilística. La previsión
puntual proporciona un único valor para cada hora, normalmente una
estimación de la media o de la mediana de la demanda futura. La
probabilística cuantifica además la incertidumbre, mediante intervalos de
predicción, cuantiles o la distribución completa de la demanda futura. Los
conceptos necesarios para evaluar ambas se presentan en las secciones de
evaluación de previsiones y de previsión probabilística de este capítulo.

% =====================================================================
\section{Fundamentos de series temporales}\label{sec:fundamentos}

Esta sección resume los conceptos de la teoría clásica de series temporales
que sostienen el modelo SARIMAX y los contrastes del análisis exploratorio.
Las referencias generales son los manuales de \textcite{box2015time},
\textcite{shumway2017time} y \textcite{hyndman2021fpp3}.

La sección comienza con los conceptos de proceso estocástico y
estacionariedad y con las funciones de autocorrelación simple y parcial, que
permiten describir la dependencia temporal de una serie. A continuación
presenta los modelos AR, MA y ARMA, el tratamiento de la tendencia mediante
integración y diferenciación y los contrastes de raíz unitaria y de
estacionariedad, y termina con la metodología Box-Jenkins, que articula estos
elementos en un procedimiento de identificación, estimación y diagnóstico.

\subsection{Procesos estocásticos y estacionariedad}\label{sec:estacionariedad}

Una serie temporal $y_1, y_2, \dots, y_T$ se interpreta como una realización
de un proceso estocástico, una sucesión de variables aleatorias indexadas
por el tiempo. Para estudiar el proceso a partir de una única realización es
necesario suponer que algunas de sus propiedades no cambian con el tiempo.
Un proceso es estacionario en sentido débil si su media es constante, su
varianza es finita y la covarianza entre dos observaciones depende solo de
la distancia entre ellas y no del momento en que se miden:
\begin{equation}
	\mathrm{E}(y_t) = \mu, \qquad
	\mathrm{Cov}(y_t, y_{t-k}) = \gamma_k \quad \text{para todo } t.
\end{equation}
La función $\gamma_k$ es la función de autocovarianza, y $\gamma_0$ es la
varianza del proceso. El caso más sencillo de proceso estacionario es el
ruido blanco, una sucesión $\varepsilon_t$ de variables con media cero,
varianza constante $\sigma^2$ y sin correlación entre sí. El ruido blanco no
contiene ninguna regularidad aprovechable para la previsión, y por eso un
modelo se considera adecuado cuando sus residuos se comportan como tal.

Una serie de demanda eléctrica no es estacionaria en sentido estricto, porque
su media cambia con la hora, el día de la semana y la estación. La estrategia
habitual consiste en eliminar o modelizar esas fuentes de no estacionariedad,
ya sea mediante componentes deterministas, como términos estacionales o de
tendencia, o mediante diferenciación, y tratar como estacionario lo que
queda.

\subsection{Funciones de autocorrelación simple y parcial}\label{sec:acf}

La función de autocorrelación (ACF) normaliza la autocovarianza,
$\rho_k = \gamma_k / \gamma_0$, y mide la correlación lineal entre
observaciones separadas $k$ periodos. Se estima a partir de la serie como
\begin{equation}
	r_k = \frac{\sum_{t=k+1}^{T} (y_t - \bar{y})(y_{t-k} - \bar{y})}
	           {\sum_{t=1}^{T} (y_t - \bar{y})^2}.
\end{equation}
La función de autocorrelación parcial (PACF), $\phi_{kk}$, mide la
correlación entre $y_t$ e $y_{t-k}$ una vez descontado el efecto lineal de
las observaciones intermedias $y_{t-1}, \dots, y_{t-k+1}$. Equivale al último
coeficiente de una regresión de $y_t$ sobre sus $k$ primeros retardos, y
permite distinguir las relaciones directas de las que se transmiten a través
de los retardos intermedios.

Si la serie fuera ruido blanco, los coeficientes estimados se distribuirían
aproximadamente como una normal de media cero y desviación típica
$1/\sqrt{T}$, de modo que la banda $\pm 1{,}96/\sqrt{T}$ delimita los valores
compatibles con la ausencia de correlación al 95\,\%. Con series largas, esa
banda es muy estrecha: con $T = 52\,583$ observaciones, como en el periodo de
entrenamiento de este trabajo, es de $\pm 0{,}0085$, y prácticamente
cualquier correlación resulta significativa. En ese contexto, el valor de
las funciones interesa más por su forma y su magnitud que por su
significación.

\subsection{Modelos AR, MA y ARMA}\label{sec:arma}

Un proceso autorregresivo de orden $p$, AR($p$), expresa cada observación
como una combinación lineal de las $p$ anteriores más una perturbación de
ruido blanco:
\begin{equation}
	y_t = c + \phi_1 y_{t-1} + \dots + \phi_p y_{t-p} + \varepsilon_t.
\end{equation}
Un proceso de medias móviles de orden $q$, MA($q$), la expresa como una
combinación de la perturbación actual y de las $q$ anteriores:
\begin{equation}
	y_t = \mu + \varepsilon_t + \theta_1 \varepsilon_{t-1} + \dots + \theta_q \varepsilon_{t-q}.
\end{equation}
El modelo ARMA($p,q$) combina ambos. Con el operador de retardo $B$, definido
por $B y_t = y_{t-1}$, se escribe de forma compacta como
\begin{equation}
	\phi(B)\, y_t = c + \theta(B)\, \varepsilon_t,
\end{equation}
donde $\phi(B) = 1 - \phi_1 B - \dots - \phi_p B^p$ y
$\theta(B) = 1 + \theta_1 B + \dots + \theta_q B^q$ son polinomios en $B$.
El proceso es estacionario si las raíces de $\phi(z) = 0$ están fuera del
círculo unidad, e invertible, es decir, expresable como un autorregresivo de
orden infinito, si lo están las de $\theta(z) = 0$.

Cada tipo de proceso deja una huella característica en las funciones de
autocorrelación. En un AR($p$), la ACF decrece de forma gradual, exponencial
u oscilante, y la PACF se anula a partir del retardo $p$. En un MA($q$)
ocurre lo contrario: la ACF se anula a partir del retardo $q$ y la PACF
decrece gradualmente. En un ARMA, ambas decrecen sin cortarse. Estos patrones
son la base de la fase de identificación de la metodología Box-Jenkins
(sección~\ref{sec:box-jenkins}).

\subsection{Tendencia, integración y diferenciación}\label{sec:integracion}

Una serie puede no ser estacionaria por dos motivos de naturaleza distinta.
En el primero, la serie fluctúa de forma estacionaria alrededor de una
componente determinista, como una tendencia lineal,
$y_t = \alpha + \beta t + u_t$ con $u_t$ estacionario: las perturbaciones
tienen un efecto transitorio y la serie tiende a volver a la tendencia. En
el segundo, la no estacionariedad es estocástica. El ejemplo más sencillo es
el paseo aleatorio, $y_t = y_{t-1} + \varepsilon_t$, en el que cada
perturbación se acumula y tiene un efecto permanente. Se dice que este
proceso tiene una raíz unitaria, porque su polinomio autorregresivo,
$1 - B$, se anula en $z = 1$.

Un proceso con raíz unitaria se vuelve estacionario al diferenciarlo,
$\nabla y_t = (1 - B)\, y_t = y_t - y_{t-1}$. Una serie que necesita $d$
diferencias para ser estacionaria se denomina integrada de orden $d$, y el
modelo ARIMA($p,d,q$) aplica un ARMA a la serie diferenciada:
\begin{equation}
	\phi(B)\,(1 - B)^d\, y_t = c + \theta(B)\, \varepsilon_t.
\end{equation}
La distinción entre los dos tipos de no estacionariedad tiene consecuencias
prácticas. Diferenciar una serie que solo tiene una tendencia determinista
introduce en ella una estructura de medias móviles no invertible, lo que se
conoce como sobrediferenciación; no diferenciar una serie con raíz unitaria
conduce a modelos con propiedades estadísticas engañosas. De ahí la
necesidad de contrastes formales.

\subsection{Contrastes de raíz unitaria y de estacionariedad}\label{sec:raiz-unitaria}

El contraste de Dickey-Fuller aumentado (ADF) parte de la regresión
\begin{equation}
	\nabla y_t = \alpha + \gamma\, y_{t-1} + \sum_{i=1}^{m} \delta_i\, \nabla y_{t-i} + \varepsilon_t,
\end{equation}
y contrasta la hipótesis nula de raíz unitaria, $\gamma = 0$, frente a la
alternativa de estacionariedad, $\gamma < 0$
\parencite{dickey1979distribution}. Los retardos de la diferencia absorben la
autocorrelación de la serie, de modo que el contraste es válido para
procesos ARMA de orden desconocido si $m$ crece con el tamaño de la muestra
\parencite{said1984testing}; en la práctica, $m$ se elige con un criterio de
información. Bajo la hipótesis nula, el estadístico del contraste no sigue
la distribución $t$ habitual, sino una distribución propia tabulada por
simulación. La regresión puede incluir solo una constante, como en la
expresión anterior, o también una tendencia lineal.

El contraste KPSS invierte las hipótesis
\parencite{kwiatkowski1992testing}. Descompone la serie en una tendencia
determinista, un paseo aleatorio y un error estacionario, y contrasta la
hipótesis nula de que la varianza del paseo aleatorio es nula, es decir, de
que la serie es estacionaria en nivel o alrededor de una tendencia, frente a
la alternativa de raíz unitaria.

Como sus hipótesis nulas son opuestas, ambos contrastes suelen emplearse
juntos. Si el ADF rechaza y el KPSS no, la evidencia apunta a una serie
estacionaria; si ocurre lo contrario, a una raíz unitaria. Si ambos
rechazan, la serie no se ajusta bien a ninguna de las dos hipótesis, lo que
es habitual cuando presenta componentes deterministas que cambian en el
tiempo, como estacionalidades o roturas de nivel. Conviene tener presente,
además, que con muestras muy grandes ambos contrastes tienen una potencia
tan alta que rechazan ante desviaciones mínimas, por lo que su resultado
debe interpretarse junto con el análisis gráfico.

La estacionalidad plantea una cuestión análoga. Puede ser determinista, un
patrón fijo que se repite con un ruido estacionario alrededor, o
estocástica, como en el paseo aleatorio estacional
$y_t = y_{t-s} + \varepsilon_t$, en el que el patrón de cada ciclo hereda el
del anterior con perturbaciones que se acumulan. En el segundo caso, la
serie tiene raíces unitarias en las frecuencias estacionales y requiere una
diferencia estacional, $(1 - B^s)\, y_t$. Existen contrastes específicos para
esta distinción, como el de \textcite{hylleberg1990seasonal}, que contrasta
raíces unitarias en cada frecuencia estacional, o el de
\textcite{canova1995seasonal}, que toma como hipótesis nula una
estacionalidad determinista estable. Ambos se concibieron para periodos
estacionales cortos, como los de las series trimestrales o mensuales, y su
aplicación a periodos de 168 u 8766 horas, con decenas o miles de
frecuencias estacionales, resulta poco práctica. Los
contrastes ADF y KPSS se refieren a la raíz unitaria ordinaria y no
resuelven esta cuestión.

\subsection{Metodología Box-Jenkins}\label{sec:box-jenkins}

La metodología de \textcite{box2015time} organiza la construcción de un
modelo ARIMA en un ciclo de tres fases que se repite hasta obtener un modelo
adecuado.

En la \emph{identificación}, se decide si la serie necesita transformaciones
o diferencias y se proponen órdenes tentativos $p$ y $q$ a partir de la
forma de la ACF y la PACF. En la \emph{estimación}, los parámetros se
obtienen por máxima verosimilitud, que en las implementaciones habituales se
calcula representando el modelo en forma de espacio de estados y aplicando
el filtro de Kalman. Cuando hay varios modelos candidatos, se comparan con
criterios de información que penalizan la complejidad, como el de Akaike,
\begin{equation}
	\mathrm{AIC} = -2 \log L + 2k,
\end{equation}
donde $L$ es la verosimilitud maximizada y $k$ el número de parámetros
\parencite{akaike1974new}; se prefiere el modelo con menor valor.

En el \emph{diagnóstico}, se comprueba que los residuos se comportan como
ruido blanco. Además de la inspección de su ACF, se emplea el contraste de
\textcite{ljung1978measure}, que resume en un único estadístico las
autocorrelaciones de los residuos hasta el retardo $h$:
\begin{equation}
	Q = T(T+2) \sum_{k=1}^{h} \frac{r_k^2}{T-k}.
\end{equation}
Bajo la hipótesis nula de ausencia de autocorrelación, $Q$ sigue
aproximadamente una distribución $\chi^2$ con $h$ menos el número de
parámetros ARMA estimados como grados de libertad. Un rechazo indica que los
residuos conservan estructura aprovechable y que el modelo debe revisarse.

Una vez aceptado el modelo, la previsión a $h$ pasos es la esperanza
condicionada de $y_{T+h}$ dada la información disponible. Si las
perturbaciones son normales, el modelo proporciona también la varianza del
error de previsión, que crece con el horizonte, y con ella intervalos de
predicción analíticos de la forma $\hat{y}_{T+h} \pm z\, \hat{\sigma}_h$.
Estos intervalos suponen que el modelo está bien especificado, que los
errores son normales y que los parámetros son conocidos, de modo que tienden
a ser demasiado estrechos cuando alguna de esas condiciones falla. Por eso su
calibración debe comprobarse empíricamente.

% =====================================================================
\section{Estacionalidad múltiple}\label{sec:estacionalidad-multiple}

La demanda horaria superpone tres ciclos, de 24, 168 y unas 8766 horas, y los
modelos ARIMA estacionales solo admiten uno. Esta sección presenta las
herramientas para detectar y describir varios ciclos simultáneos, las
limitaciones del modelo SARIMA y las dos familias de modelos que las
superan: la regresión con términos de Fourier y errores ARIMA, que es la base
del SARIMAX de este trabajo, y los modelos de suavizado exponencial con
estacionalidad compleja.

\subsection{Análisis espectral y periodograma}\label{sec:espectral}

El análisis espectral parte de una idea sencilla: cualquier serie de longitud
$N$ puede escribirse como suma de ondas sinusoidales de frecuencias
$\omega_k = k/N$, con $k = 1, \dots, N/2$, llamadas frecuencias de Fourier
\parencite{shumway2017time}. La frecuencia se mide en ciclos por unidad de
tiempo, y su inversa es el periodo: una frecuencia de $1/24$ ciclos por hora
corresponde a un ciclo de 24 horas. El periodograma mide cuánto contribuye
cada frecuencia a la serie:
\begin{equation}
	I(\omega_k) = \frac{1}{N} \left| \sum_{t=1}^{N} (y_t - \bar{y})\, e^{-2\pi i \omega_k t} \right|^2.
\end{equation}
Su valor en cada frecuencia es proporcional al cuadrado de la amplitud de la
onda correspondiente, y la suma de todos sus valores es proporcional a la
varianza de la serie. El periodograma reparte, por tanto, la varianza entre
ciclos de distinta duración, y los ciclos importantes aparecen como picos.

Su lectura requiere tener en cuenta tres propiedades. La primera es la
resolución. Las frecuencias de Fourier están igualmente espaciadas, a
distancia $1/N$, pero los periodos correspondientes no: cerca de un periodo
$P$, los periodos contiguos están aproximadamente a $P^2/N$ de distancia, de
modo que la resolución empeora para los periodos largos. Con unas 52\,600
observaciones horarias, por ejemplo, los periodos de la rejilla próximos al
año son de unas 7500, 8760 y 10\,500 horas, y el periodograma no puede
distinguir ciclos cercanos al anual.

La segunda son los armónicos. Un ciclo de periodo $P$ que no tiene forma
sinusoidal aparece en el periodograma no solo en $P$, sino también en sus
divisores $P/2$, $P/3$, etc. Un perfil diario con dos máximos, por ejemplo,
genera un pico en 12 horas además de en 24. Es la misma idea que permite
representar cualquier patrón periódico con términos de Fourier
(sección~\ref{sec:fourier}): cuanto más se aleja su forma de una sinusoide,
más armónicos hacen falta.

La tercera es la fuga espectral. Cuando una componente de la serie no
coincide exactamente con una frecuencia de Fourier, o cuando la serie tiene
tendencia o cambios de nivel, su potencia se reparte entre las frecuencias
vecinas, sobre todo hacia las más bajas. Existen técnicas de suavizado y de
ventaneo que reducen este efecto y la elevada variabilidad del periodograma
bruto, que no disminuye al aumentar la longitud de la serie
\parencite{shumway2017time}; para localizar los picos dominantes, sin
embargo, el periodograma bruto suele ser suficiente.

\subsection{Descomposición STL y MSTL}\label{sec:mstl}

Las técnicas de descomposición separan una serie en componentes
interpretables. En la descomposición aditiva, la serie es la suma de una
tendencia $T_t$, una componente estacional $S_t$ y un residuo $R_t$:
$y_t = T_t + S_t + R_t$. El método STL \parencite{cleveland1990stl} realiza
esta separación mediante regresiones locales ponderadas, denominadas LOESS,
aplicadas de forma iterativa a la tendencia y a cada posición del ciclo
estacional. Tiene dos propiedades que lo hacen adecuado para series
largas: la componente estacional puede evolucionar lentamente a lo largo de
la serie, con una velocidad que regula el tamaño de la ventana de suavizado,
y el procedimiento admite una versión robusta que reduce la influencia de los
valores atípicos.

STL solo admite un periodo estacional. El método MSTL
\parencite{bandara2025mstl} lo extiende a varios periodos aplicando STL de
forma iterativa a cada uno de ellos, empezando por el más corto, y
refinando las estimaciones en varias pasadas. El resultado es una
componente estacional por periodo, además de la tendencia y el residuo.

En este trabajo, la descomposición es una herramienta descriptiva: permite
ver cómo evoluciona cada ciclo y qué parte de la variabilidad no explican los
ciclos regulares. Tiene, no obstante, limitaciones que condicionan su
lectura. Las estimaciones son menos fiables en los extremos de la serie,
donde el suavizado dispone de datos solo por un lado; una componente con
pocos ciclos completos, como la anual en una serie de pocos años, queda poco
determinada; y la descomposición aditiva supone que la amplitud de los ciclos
no depende del nivel de la serie.

\subsection{El modelo SARIMA y su limitación}\label{sec:sarima}

El modelo ARIMA estacional, SARIMA$(p,d,q)(P,D,Q)_s$, incorpora la
estacionalidad de periodo $s$ mediante polinomios en $B^s$
\parencite{box2015time}:
\begin{equation}
	\phi(B)\, \Phi(B^s)\, (1 - B)^d\, (1 - B^s)^D\, y_t = c + \theta(B)\, \Theta(B^s)\, \varepsilon_t,
\end{equation}
donde $\Phi$ y $\Theta$ son polinomios autorregresivo y de medias móviles
estacionales, de órdenes $P$ y $Q$, y $(1 - B^s)^D$ es la diferencia
estacional. El modelo relaciona cada observación con las del mismo momento
de los ciclos anteriores y es muy eficaz con series mensuales o trimestrales.

Su limitación fundamental para la demanda horaria es que admite un único
periodo $s$. Con $s = 24$ recoge el ciclo diario, pero no el semanal ni el
anual. Con $s = 168$, los polinomios estacionales relacionan cada hora con la
misma hora de las semanas anteriores, pero la estimación con periodos tan
largos resulta poco práctica \parencite{hyndman2021fpp3}; con el periodo
anual, de 8766 horas, es inviable, y además ese periodo no es entero en
días, porque un año medio tiene 365,25. Encadenar varias diferencias
estacionales tampoco resuelve el problema: cada una supone que la
estacionalidad correspondiente es estocástica y reduce la información
disponible.

\subsection{Regresión con errores ARIMA y términos de Fourier}\label{sec:fourier}

Una forma de superar la limitación anterior es separar la estacionalidad de
la dinámica de corto plazo. La regresión con errores ARIMA modela la serie
como una regresión lineal sobre un conjunto de variables explicativas
$\mathbf{x}_t$, con un error que sigue un proceso ARIMA en lugar de ser ruido
blanco \parencite{hyndman2021fpp3}:
\begin{equation}
	y_t = \boldsymbol{\beta}^{\top} \mathbf{x}_t + \eta_t,
	\qquad \phi(B)\, (1-B)^d\, \eta_t = \theta(B)\, \varepsilon_t.
\end{equation}
Las variables explicativas recogen los efectos deterministas y externos,
como la estacionalidad, el calendario o la temperatura, y el proceso ARIMA
de los errores recoge la dependencia de corto plazo que queda tras ellos. El
modelo SARIMAX de la librería statsmodels, empleado en este trabajo, se
formula de esta manera, como una regresión con errores SARIMA.

Para representar la estacionalidad se emplean términos de Fourier. Para un
periodo $m$, se definen $K$ parejas de variables
\begin{equation}
	s_k(t) = \sin\!\left(\frac{2\pi k t}{m}\right),
	\qquad
	c_k(t) = \cos\!\left(\frac{2\pi k t}{m}\right),
	\qquad k = 1, \dots, K,
\end{equation}
con $K \le m/2$. Una combinación lineal de estas variables puede reproducir
cualquier patrón periódico de periodo $m$: con $K = m/2$ la representación es
exacta, y con valores menores se obtiene una aproximación suave, tanto más
fiel cuanto mayor es $K$. Los términos de Fourier tienen tres ventajas sobre
el SARIMA. Admiten varios periodos a la vez, basta con incluir los términos
de cada uno; admiten periodos no enteros, como el anual; y necesitan pocos
parámetros, $2K$ por periodo, frente a los $m - 1$ de un conjunto de
variables indicadoras. El modelo que combina términos de Fourier con errores
ARIMA se conoce como regresión armónica dinámica
\parencite{hyndman2021fpp3}.

El número de armónicos $K$ de cada periodo controla el compromiso entre
flexibilidad y sobreajuste, y se elige con un criterio de información o con
el error en un periodo de validación. Su principal limitación es que el
patrón estacional es fijo: los coeficientes de los términos de Fourier son
constantes, de modo que el modelo impone la misma forma del ciclo diario en
enero que en julio. Cuando el perfil cambia a lo largo del año, solo puede
recogerse añadiendo interacciones, por ejemplo entre los términos diarios y
la temperatura o la estación, a costa de multiplicar el número de
parámetros.

\subsection{Suavizado exponencial y TBATS}\label{sec:tbats}

La otra gran familia de modelos clásicos es la del suavizado exponencial,
que construye la previsión como una media ponderada de las observaciones
pasadas con pesos que decrecen exponencialmente con la antigüedad. En su
formulación moderna, como modelos de espacio de estados con componentes de
error, tendencia y estacionalidad, proporciona también intervalos de
predicción \parencite{hyndman2021fpp3}. El método de Holt-Winters, su versión
estacional más conocida, admite un único periodo. \textcite{taylor2003short}
lo extendió a dos estacionalidades, la diaria y la semanal, precisamente para
la previsión de la demanda eléctrica a corto plazo.

El modelo TBATS \parencite{delivera2011forecasting} generaliza esta idea a un
número arbitrario de periodos, incluidos los no enteros. Combina una
transformación de Box-Cox, errores ARMA, una tendencia y componentes
estacionales trigonométricas, análogas a los términos de Fourier, pero cuyos
coeficientes son estados del modelo y evolucionan en el tiempo. Esa es su
principal ventaja frente a la regresión armónica: la forma de cada ciclo
puede cambiar lentamente a lo largo de la serie, en lugar de repetirse
invariable. Su inconveniente es el coste computacional, que crece con la
longitud de la serie y con el número y la duración de los periodos
estacionales y que, en una serie horaria de varios años con un periodo
anual, puede ser elevado. Las consecuencias de este coste para el diseño del
trabajo se discuten en el marco metodológico.
